% GENERATED FILE. Edit canonical lessons, then run npm run generate:books.
% canonical-source-hash: fnv1a64:5c5fb21397cadc8e
% canonical-lessons: ES-C571-inventor, ES-C571-descubrimiento, ES-C571-invento, ES-C571-origen, ES-C571-desarrollo

\chapter{Un gran descubrimiento}
\label{ch:es-un-gran-descubrimiento}
\hlchaptermodality{571}

\begin{chapteropening}
\textbf{By the end of this chapter:} \emph{I can talk about how things began: an inventor, a discovery, an invention, its origin and how it developed.}

Complete the last lesson of chapter 571: I can talk about how things began: an inventor, a discovery, an invention, its origin and how it developed.
\end{chapteropening}

\section[el inventor]{el inventor — the inventor}

\label{lesson:ES-C571-inventor}



\begin{quote}
Before the new one: say the Spanish for the scandal, then the Spanish for the controversy.
\end{quote}

\subsection*{You'll want to know: el inventor}
\textbf{el inventor} — \textquotedblleft{}the inventor\textquotedblright{}.

\begin{quote}
El inventor del teléfono fue muy famoso.
\end{quote}

\begin{quote}
\emph{The inventor of the telephone was very famous.}
\end{quote}

A woman is \emph{la inventora}.

\begin{cousinweb}
\emph{inventor} is a learned borrowing from Latin \emph{inventor} 'one who finds out, a deviser' (EtymDB), from \emph{invenīre} 'to come upon' (Lewis \& Short).
\end{cousinweb}

\begin{grammarlens}[title={What you've built}]
One more word for telling, explaining and describing at B1.
\end{grammarlens}

\subsection*{Your turn}
\begin{itemize}
\raggedright
  \item \emph{Say it:} \emph{el inventor}
  \item \emph{Say it:} \emph{el inventor}, once more
  \item \emph{Say it:} a sentence of your own with \emph{el inventor}
  \item \emph{Recall:} say the Spanish for the scandal, then the Spanish for the controversy, then say \emph{el inventor} again
\end{itemize}

\begin{tcolorbox}[breakable,colback=teal!4,colframe=teal!35!black,title={Before you move on}]
What does el inventor mean? (\textquotedblleft{}The inventor\textquotedblright{}.) Say it once more.
\end{tcolorbox}
\section[el descubrimiento]{el descubrimiento — the discovery}

\label{lesson:ES-C571-descubrimiento}



\begin{quote}
Before the new one: say the Spanish for the controversy, then the Spanish for the inventor.
\end{quote}

\subsection*{You'll want to know: el descubrimiento}
\textbf{el descubrimiento} — \textquotedblleft{}the discovery\textquotedblright{}.

\begin{quote}
Fue un descubrimiento muy importante para la medicina.
\end{quote}

\begin{quote}
\emph{It was a very important discovery for medicine.}
\end{quote}

\begin{cousinweb}
\emph{descubrimiento} is built on \emph{descubrir} with the ending \emph{-miento}. \emph{Descubrir} comes from \emph{cubrir}, inherited from Latin \emph{cooperīre} 'to cover over' (Lewis \& Short).
\end{cousinweb}

\begin{grammarlens}[title={What you've built}]
One more word for telling, explaining and describing at B1.
\end{grammarlens}

\subsection*{Your turn}
\begin{itemize}
\raggedright
  \item \emph{Say it:} \emph{el descubrimiento}
  \item \emph{Say it:} \emph{el descubrimiento}, once more
  \item \emph{Say it:} a sentence of your own with \emph{el descubrimiento}
  \item \emph{Recall:} say the Spanish for the controversy, then the Spanish for the inventor, then say \emph{el descubrimiento} again
\end{itemize}

\begin{tcolorbox}[breakable,colback=teal!4,colframe=teal!35!black,title={Before you move on}]
What does el descubrimiento mean? (\textquotedblleft{}The discovery\textquotedblright{}.) Say it once more.
\end{tcolorbox}
\section[el invento]{el invento — the invention}

\label{lesson:ES-C571-invento}



\begin{quote}
Before the new one: say the Spanish for the inventor, then the Spanish for the discovery.
\end{quote}

\subsection*{You'll want to know: el invento}
\textbf{el invento} — \textquotedblleft{}the invention\textquotedblright{}.

\begin{quote}
El teléfono fue un gran invento.
\end{quote}

\begin{quote}
\emph{The telephone was a great invention.}
\end{quote}

\emph{Inventar} is to invent, and also to make something up.

\begin{cousinweb}
\emph{invento} comes from Latin \emph{inventum} 'a device, an invention' (EtymDB), from \emph{invenīre} 'to come upon, to find'.
\end{cousinweb}

\begin{grammarlens}[title={What you've built}]
One more word for telling, explaining and describing at B1.
\end{grammarlens}

\subsection*{Your turn}
\begin{itemize}
\raggedright
  \item \emph{Say it:} \emph{el invento}
  \item \emph{Say it:} \emph{el invento}, once more
  \item \emph{Say it:} a sentence of your own with \emph{el invento}
  \item \emph{Recall:} say the Spanish for the inventor, then the Spanish for the discovery, then say \emph{el invento} again
\end{itemize}

\begin{tcolorbox}[breakable,colback=teal!4,colframe=teal!35!black,title={Before you move on}]
What does el invento mean? (\textquotedblleft{}The invention\textquotedblright{}.) Say it once more.
\end{tcolorbox}
\section[el origen]{el origen — the origin}

\label{lesson:ES-C571-origen}



\begin{quote}
Before the new one: say the Spanish for the discovery, then the Spanish for the invention.
\end{quote}

\subsection*{You'll want to know: el origen}
\textbf{el origen} — \textquotedblleft{}the origin\textquotedblright{}.

\begin{quote}
Nadie conoce el origen del incendio.
\end{quote}

\begin{quote}
\emph{Nobody knows the origin of the fire.}
\end{quote}

Plural: \emph{los orígenes}.

\begin{cousinweb}
\emph{origen} is a learned borrowing from Latin \emph{orīgō} 'a beginning, a source' (EtymDB), from \emph{orīrī} 'to rise'.
\end{cousinweb}

\begin{grammarlens}[title={What you've built}]
One more word for telling, explaining and describing at B1.
\end{grammarlens}

\subsection*{Your turn}
\begin{itemize}
\raggedright
  \item \emph{Say it:} \emph{el origen}
  \item \emph{Say it:} \emph{el origen}, once more
  \item \emph{Say it:} a sentence of your own with \emph{el origen}
  \item \emph{Recall:} say the Spanish for the discovery, then the Spanish for the invention, then say \emph{el origen} again
\end{itemize}

\begin{tcolorbox}[breakable,colback=teal!4,colframe=teal!35!black,title={Before you move on}]
What does el origen mean? (\textquotedblleft{}The origin\textquotedblright{}.) Say it once more.
\end{tcolorbox}
\section[el desarrollo]{el desarrollo — the development}

\label{lesson:ES-C571-desarrollo}



\begin{quote}
Before the new one: say the Spanish for the invention, then the Spanish for the origin.
\end{quote}

\subsection*{You'll want to know: el desarrollo}
\textbf{el desarrollo} — \textquotedblleft{}the development\textquotedblright{}.

\begin{quote}
El desarrollo de la ciudad fue muy rápido.
\end{quote}

\begin{quote}
\emph{The city's development was very fast.}
\end{quote}

\begin{cousinweb}
\emph{desarrollo} is built on \emph{desarrollar} 'to unroll, to develop', from \emph{des-} and \emph{arrollar} 'to roll up', built on \emph{rollo} 'a roll', which is usually traced to Latin \emph{rotulus} 'a little wheel'.
\end{cousinweb}

\begin{grammarlens}[title={What you've built}]
One more word for telling, explaining and describing at B1.
\end{grammarlens}

\subsection*{Your turn}
\begin{itemize}
\raggedright
  \item \emph{Say it:} \emph{el desarrollo}
  \item \emph{Say it:} \emph{el desarrollo}, once more
  \item \emph{Say it:} a sentence of your own with \emph{el desarrollo}
  \item \emph{Recall:} say the Spanish for the invention, then the Spanish for the origin, then say \emph{el desarrollo} again
\end{itemize}

\begin{tcolorbox}[breakable,colback=teal!4,colframe=teal!35!black,title={Before you move on}]
What does el desarrollo mean? (\textquotedblleft{}The development\textquotedblright{}.) Say it once more.
\end{tcolorbox}
